\subsection{关于点测度积分的正函数上积分}

我们将看到，当 \((\pi,g)\) 是一个 \(\mu\) -适应对时，把 §3 的诸命题应用于族 \(t\mapsto\lambda_{t}=g(t)\varepsilon_{\pi(t)}\) （由命题 1，它是 \(\mu\) -适宜的）所得到的结果可以被加强。

定理 1. --- 设 \((\pi, g)\) 是一个 \(\mu\) -适应对，并令

\[
\nu = \int g (t) \varepsilon_ {\pi (t)} d \mu (t).
\]

对每个数值函数 \(f \geqslant 0\) （定义在 \(X\) 上），

\[
\int^ {\bullet} f (x) d \nu (x) = \int^ {\bullet} f (\pi (t)) g (t) d \mu (t).\tag{1}
\]

\begin{enumerate}
\def\labelenumi{\Alph{enumi})}
\tightlist
\item
  首先假设测度 \(\mu\) 具有紧支集 \(K\) ，且限制到 \(K\) 上时函数 \(g\) 和 \(\pi\) 都连续。由 §3, No.~1 的公式 (4)，\(\nu^{\bullet}(1) = \int_{K} g(t) d\mu(t) < +\infty\) ，故出现在公式 (1) 中的所有测度都是有界的。因此我们可以在第一成员中把 \(\int^{\bullet}\) 换成 \(\int^{*}\) 。鉴于 §3, No.~2 的公式 (6)，问题归结为证明
\end{enumerate}

\[
\int^ {*} f (x) d \nu (x) \leqslant \int^ {\bullet} f (\pi (t)) g (t) d \mu (t),\tag{2}
\]

其中第二成员中的记号 \(\int^{\bullet}\) 又可以换成 \(\int^{*}\) 。由上积分的定义，只需验证不等式

\[
\int^ {*} f (x) d \nu (x) \leqslant \int^ {*} h (t) d \mu (t)\tag{3}
\]

对每个下半连续函数 \(h\) （\(\mathrm{T}\) 上的）——它 \(\geqslant\) 函数 \(t \mapsto f(\pi(t))g(t)\) 。现在，设 \(\varepsilon\) 是一个数 \(> 0\) ，\(u\) 是函数 \((h + \varepsilon)/g\) ，它在 \(\mathrm{K}\) 中是下半连续的。如果 \(t \in \overline{\pi}^{-1}(\{x\}) \cap \mathrm{K}\) ，那么 \(u(t) \geqslant f(x)\) ：当 \(g(t) = 0\) 时这是明显的，因为那时 \(u(t) = +\infty\) ；如果 \(g(t) > 0\) ，那么

\[
u (t) g (t) = h (t) + \varepsilon \geqslant f (\pi (t)) g (t) = f (x) g (t),
\]

由此得出所要的不等式。在这些条件下，对每个 \(x \in X\) ，令 \(v(x)\) 为 \(u(t)\) 对 \(t \in \overline{\pi}^{-1}(\{x\}) \cap K\) 的下确界。由前所述，函数 v \(\geqslant f\) ；由引理（应用于 \(\pi\) 在 K 上的限制），它在 X 上是下半连续的；且 \(v(\pi(t))g(t) \leqslant h(t) + \varepsilon\) 对所有 \(t \in K\) 成立（回忆：按约定，若 \(g(t) = 0\) ，则第一成员为零）。于是对 v 应用 §3, No.~1 的公式 (4)。我们得到：

\[
\begin{array}{l} \int^ {*} f (x) d \nu (x) \leqslant \int^ {*} v (x) d \nu (x) \\ = \int^ {*} v (\pi (t)) g (t) d \mu (t) \leqslant \int_ {K} ^ {*} (h (t) + \varepsilon) d \mu (t) \\ = \int^ {*} h (t) d \mu (t) + \varepsilon \mu (1). \end{array}\tag{4}
\]

由于测度 \(\mu\) 有界且 \(\varepsilon\) 是任意的，就得出不等式 (3)。

\begin{enumerate}
\def\labelenumi{\Alph{enumi})}
\setcounter{enumi}{1}
\tightlist
\item
  现在转到一般情形。由于映射 \(t \mapsto (\pi(t), g(t))\) （\(T\) 到 \(X \times R_{+}\) 中的）是 \(\mu\) -可测的（Ch. IV, §5, No.~3, 定理 1），集 \(\mathfrak{K}\) （由紧子集 \(K\) （\(T\) 中的）组成，使 \(\pi\) 和 \(g\) 到 \(K\) 上的限制都连续）是 \(\mu\) -稠密的（Ch. IV, §5, No.~10, 命题 15）。由 §2, No.~3 的命题 4，\(\mu\) 是一个可和族 \((\mu_{\alpha})_{\alpha \in A}\) 之和，这些测度的支集都是 \(\mathfrak{K}\) 的元素；对 \((\pi, g)\) 是 \(\mu_{\alpha}\) -适应的（对每个 \(\alpha \in A\) ），令 \(\nu_{\alpha}\) 为测度 \(\int g(t)\varepsilon_{\pi(t)} d\mu_{\alpha}(t)\) 。于是由 A)，
\end{enumerate}

\[
\int^ {\bullet} f (x) d \nu_ {\alpha} (x) = \int^ {\bullet} f (\pi (t)) g (t) d \mu_ {\alpha} (t).\tag{5}
\]

但诸 \(\nu_{\alpha}\) 构成一个可和族，其和等于 \(\nu\) （§3, No.~1, 命题 1 的推论）。因此，由 §2, No.~2 的命题 1，

\[
\int^ {\bullet} f (x) d \nu (x) = \sum_ {\alpha \in \mathrm{A}} \int^ {\bullet} f (x) d \nu_ {\alpha} (x).\tag{6}
\]

对 (5) 的第二成员有类似的关系，因此对 \(\alpha\) 求和便从 (5) 得出 (1)。

推论. --- X 的子集 N 为局部 \(\nu\) -可忽略的，必须且只须 \(\overline{\pi}^{-1}(N)\) 与点 \(t \in T\) （满足 \(g(t) > 0\) 者）所成的集的交是局部 \(\mu\) -可忽略的。

命题 2. --- 设 \(\pi\) 是 T 到 X 中的一个逆紧连续映射（GT, I, §10, No.~1），\(g\) 是 T 上的一个连续有限数值函数，使 \(g(t) > 0\) 对所有 \(t \in \mathrm{T}\) 成立。则对 \((\pi, g)\) 是 \(\mu\) -适应的，且若令

\[
\nu = \int g (t) \varepsilon_ {\pi (t)} d \mu (t),
\]

那么对每个数值函数 \(f \geqslant 0\) （定义在 \(X\) 上），

\[
\int^ {*} f (x) d \nu (x) = \int^ {*} f (\pi (t)) g (t) d \mu (t).\tag{7}
\]

显然 \(\pi\) 和 g 都是 \(\mu\) -可测的；此外，对每个函数 \(\psi \in \mathcal{K}(\mathrm{X})\) ，\(\psi \circ \pi\) 都是具有紧支集的连续函数，因为 \(\pi\) 是逆紧的；因此对 \((\pi, g)\) 是 \(\mu\) -适应的，而且映射 \(t \mapsto g(t) \varepsilon_{\pi(t)}\) 是淡连续的。

设 h 是 T 上的一个下半连续函数，使得

\[
f (\pi (t)) g (t) \leqslant h (t) \quad \text { for   all } t \in \mathrm{T}.
\]

我们将证明

\[
\int^ {*} f (x) d \nu (x) \leqslant \int^ {*} h (t) d \mu (t).\tag{8}
\]

由上积分的定义，这将蕴涵不等式

\[
\int^ {*} f (x) d \nu (x) \leqslant \int^ {*} f (\pi (t)) g (t) d \mu (t),
\]

把它与 §3, No.~2 的不等式 (7) 相结合，便证明了 (7)。

为证明 (8)，让我们定义一个函数 \(\overline{f}\) （\(X\) 上的）如下：\(\overline{f}(x)\) 是 \(h(t) / g(t)\) 在集 \(\overline{\pi}^{-1}(x)\) 中的下确界（下确界等于 \(+\infty\) ，如果 \(\overline{\pi}^{-1}(x) = \varnothing\) ）。函数 \(\overline{f}\) 具有下列性质：

\(1^{\circ}\overline{f} (x)\geqslant f(x)\) 对所有 \(x\in \mathrm{X}\) 成立（因为 \(g(t) > 0\) 对所有 \(t\in \mathrm{T}\) 成立）。

\[
2 ^ {\circ} \overline {{f}} (\pi (t)) g (t) \leqslant h (t) \text {   for   all   } t \in \mathrm{T}.
\]

\(3^{\circ}\) 函数 \(\overline{f}\) 是下半连续的，依据引理，因为函数 \(h / g\) 在 T 上是下半连续的。

因此，鉴于 §3, No.~1 的命题 2a)：

\[
\int^ {*} f (x) d \nu (x) \leqslant \int^ {*} \overline {{{f}}} (x) d \nu (x) = \int^ {*} \overline {{{f}}} (\pi (t)) g (t) d \mu (t) \leqslant \int^ {*} h (t) d \mu (t),
\]

这就建立了 (8)，并完成了证明。

